\section[Conclusão]{Conclusão}

A AGES IV foi a etapa em que mais precisei sair da minha zona de conforto técnica e assumir uma postura mais ampla de organização, acompanhamento e responsabilidade pelo andamento do projeto. Desde a Sprint 0, quando participei da criação das \textit{user stories}, da organização das tarefas do painel administrativo e da preparação da reunião com os clientes, percebi que o desafio principal não seria somente técnico. Era necessário transformar necessidades ainda abertas em tarefas claras, alinhar expectativas e criar uma base para que o time conseguisse trabalhar.

Ao longo do projeto, minha atuação ficou muito ligada à tentativa de manter o time organizado. Trabalhei no refinamento de tarefas de \textit{frontend} e \textit{backend}, acompanhei pontos ligados a edificações, imagens e idiomas, validei a modelagem do banco de dados com outros AGES IV e ajudei a reorganizar o time em \textit{squads}. Essa mudança foi importante porque melhorou a comunicação assíncrona e deu referências mais claras para quem precisava de ajuda. Mesmo assim, também ficou evidente que organização de processo não resolve tudo quando parte do time não se comunica bem ou não demonstra o mesmo nível de interesse.

Uma das maiores dificuldades da AGES IV foi lidar com ausência de respostas, falta de transparência e baixo comprometimento de alguns colegas. Em diferentes momentos, principalmente nas Sprints 1, 3 e 4, precisei cobrar retornos, tentar entender o estado das tarefas e, em alguns casos, assumir atividades que não deveriam ter chegado até mim. Isso ficou mais forte na reta final, quando precisei fazer o deploy do backend, desenvolver a autenticação com JWT, criar o CRUD de arquitetos e resolver pendências importantes para que o projeto tivesse condições de ser apresentado.

Do ponto de vista comportamental, a AGES IV foi uma experiência ambígua. Por um lado, tive oportunidades importantes de conduzir reunião com cliente, apresentar o projeto, organizar tarefas, acompanhar colegas e tomar decisões com mais maturidade. Por outro, enfrentei frustrações com comunicação, desinteresse e tarefas acumuladas no final. Eu reconheço que nosso trabalho como AGES IV poderia ter sido melhor em alguns momentos, principalmente na distribuição de responsabilidades e no acompanhamento mais firme desde o início, mas também percebi que existem limites para o quanto é possível compensar a falta de envolvimento de outras pessoas.

A parte na qual mais gostei de atuar foi na comunicação com o cliente: apresentar o projeto, sanar dúvidas e entender os requisitos sempre foram coisas que eu não me via fazendo, mas por um lado sempre quis. Desempenhar essa função foi muito bom pois eu percebi que levo jeito pra isso e que mesmo que não tivesse atuado nessa área antes desse projeto, me saí muito bem, sendo sempre objetivo e claro, mas sem perder a essência de entender as dores do cliente.

No fim, considero que a AGES IV consolidou uma visão mais realista sobre Engenharia de Software. Aprendi que um projeto não depende apenas de tecnologia, arquitetura ou boas ferramentas. Ele depende principalmente de comunicação, responsabilidade e alinhamento constante entre as pessoas. Mesmo com dificuldades, consegui contribuir para que o produto avançasse, assumi tarefas críticas quando foi necessário e desenvolvi uma percepção mais clara sobre liderança técnica, gestão de time e tomada de decisão em situações de pressão.
